%!TEX root = ../main.tex

\addtocontents{toc}{\protect\vskip1em}
\chapter*{Conclusions}\label{ch:conclusions}
\markboth{Conclusions}{}
\addcontentsline{toc}{chapter}{Conclusions}


% ===========================================================================
\section*{Summary of Findings}

This thesis presented a progressive empirical investigation of \gls{rag} for
document-grounded \gls{cqa} on the \gls{mdd} benchmark. Five \gls{rag} pipeline
configurations were evaluated on the full 4{,}427-example validation split and
on a fixed 500-example diagnostic subset, with a further two ablations on the
diagnostic subset. The key findings are as follows.

\textbf{\gls{rag} substantially improves over the \gls{llm}-only baseline.} The
\gls{llm}-only system achieves $\fone = 15.9\%$ on the full validation split.
All \gls{rag} configurations outperform this substantially, with the best system
(Hybrid+Rewrite) reaching $\fone = 26.8\%$, a gain of $+10.9$~pp. The performance
ordering is consistent across both the full split and dev-500.

\textbf{Query rewriting is the single most impactful intervention.} Adding
conversational query rewriting to hybrid retrieval raises \ratk{5} from 44.7\%
to 59.5\% on the full split ($+14.8$~pp) and improves \fone{} by $+2.0$~pp
($p = 0.017$, paired bootstrap). The improvement is driven by the rewriter
correctly anchoring retrieval to the target domain, raising domain-level
\ratk{5} from 69.4\% to 92.5\%.

\textbf{The oracle ceiling reveals a generation bottleneck.} The oracle
experiment (gold context provided directly) achieves $\fone = 36.1\%$, establishing
the generation ceiling for this model--prompt combination. The 8.1~pp gap between
the best \gls{rag} system and the oracle corresponds partly to retrieval misses
(31\% of dev-500 examples) and partly to the generation quality for retrieved
examples (F1-hit $\approx$ 33.3\%). Closing this gap requires improvements to
both generation quality (better prompts, larger models, or \gls{llm}-judge
evaluation) and retrieval coverage.

\textbf{Reranking improves retrieval but not generation.} The cross-encoder
reranker raises \ratk{5} from 61.2\% to 67.4\% ($+6.2$~pp), recovering 30
additional correct sections among the 500 dev examples. However, the
corresponding \fone{} gain of $+0.16$~pp is not statistically significant
($p = 0.414$). Once retrieval recall is sufficiently high, further improvements
yield diminishing returns because the generation quality for retrieved examples
is the binding constraint.

\textbf{Agent-style prompting reduces false abstentions without a \fone{} gain.}
The v2 prompt reduces the false-abstention rate from 7.0\% to 4.9\% on
Hybrid+Rewrite and from 4.1\% to 1.9\% on the oracle system, while abstain
precision improves. Token \fone{} is unchanged ($p > 0.3$ in all paired tests).
Prompt engineering can improve calibration but is unlikely to drive large \fone{}
gains within the current pipeline structure.


% ===========================================================================
\section*{Scope Limitations}

All ablation experiments (oracle, v2 prompt, reranker) were conducted on the
dev-500 diagnostic subset. Reported confidence intervals are correspondingly
wider than those for the full split; differences below $\sim$1~pp should be
interpreted with caution. The five base pipelines were evaluated on the full
split, but the ablation configurations (Steps 7–8 in the experimental sequence)
were not, due to the computational cost of re-running full-split experiments
with parallel \gls{llm} generation ($\sim$8\,h per pipeline).

Token \fone{} is the sole generation metric reported. As noted in
Chapter~\ref{ch:rag}, reference answers are verbatim extracts from the source
document, while the model produces paraphrases; semantically correct paraphrases
receive low token \fone{}. An \gls{llm}-based judge evaluating semantic
correctness rather than token overlap would provide a more reliable signal,
particularly for the v2 prompt ablation.


% ===========================================================================
\section*{Future Directions}

\textbf{Full-split ablations.} The most immediate next step is to run the
Hybrid+Rewrite+Rerank, Oracle (v2), and Hybrid+Rewrite (v2) pipelines on the
full validation split to validate the dev-500 findings at scale.

\textbf{LLM-judge evaluation.} Replacing or complementing token \fone{} with
a local \gls{llm} judge (e.g.\ Qwen2.5-14B-Instruct or Llama~3.1~8B) would
close the semantic gap between generation style and reference format. The
evaluation protocol would score each (question, context, prediction, reference)
tuple on a three-point scale (correct / partially correct / incorrect) and
report agreement with human judgements as a calibration measure.

\textbf{Knowledge graph augmentation.} The \gls{mdd} corpus is amenable to
\gls{kg} construction: entities such as benefit programmes, eligibility
conditions, deadlines, and fees can be extracted and linked across documents.
A hybrid pipeline that combines flat-document retrieval with graph-query
traversal (\emph{Graph~\gls{rag}}) could surface multi-hop relational evidence
that dense retrieval over raw text is likely to miss, particularly for questions
that span multiple documents or require chained inference over extracted facts.

\textbf{Agent architectures.} Beyond pipeline \gls{rag}, an agent-based
architecture could orchestrate multiple specialised tools --- a retrieval tool,
a graph-query tool, and a code-execution tool for numerical reasoning --- using
the \gls{llm} as a planner. This would address the limitation of fixed-$K$
context windows for multi-hop questions and allow the system to refine its
retrieval strategy based on intermediate findings.

\textbf{Evaluation of unanswerable calibration.} With only 31 unanswerable
examples in dev-500 (and 226 in the full split), abstention metrics are
noisy. Augmenting the dataset with additional unanswerable examples, or
evaluating on a purpose-built unanswerable-question benchmark, would enable
more reliable conclusions about abstention behaviour across pipeline variants.
